\section{Scope and limitations}
\label{sec:limitations}

The results have a deliberately narrow interpretation.  Profile~23 proves the
atomic-v3 relation in Definition~\ref{def:profile23-relation}: one input note
is replaced by one output note along the same depth-20 private path.  It does
not provide multiple inputs or outputs, a public append, recursion, proof
aggregation, a wallet protocol, a relayer, or a general private application
interface.  The pool address, sequence, old and new anchors, nullifier, output
commitment, asset identifier, and fee are public.

Theorem~\ref{thm:bcs-soundness} is classical random-oracle argument soundness
in the work-normalized metric.  It is not knowledge soundness and makes no
quantum-random-oracle claim.  Its concrete ledger includes a 124-bit
Poseidon2 assumption and treats SHA-256 as a 256-bit random oracle.  The
hiding results use the stronger programmable-random-oracle experiment; they
do not instantiate a standard-model pseudorandom-generator theorem or a
statistical zero-knowledge theorem.

The hiding view is exactly Definition~\ref{def:profile23-complete-view}.  The
equality \(\epsilon_{\rm side}=0\) is by definition of that fixed public
\(\mathsf{Proof}\)-or-\(\mathsf{Abort}\) channel.  Local filesystem access,
the burned-nonce ledger, scheduler and process timing outside the release
event, power, thermal and memory observations, and remote prover or miner
traffic are not in the modeled view.  An implementation that exposes any of
those channels requires a new simulator and leakage bound; the present
deployment model therefore keeps rejected attempts and mining local and
private.

The hiding game assumes ideal completion before the public boundary, or an
external truncation event independent of the witness and all hidden prover
state.  The evaluated controller instead enforces a concrete 480-second
deadline and publishes \(\mathsf{Abort}\) when no candidate is ready.  The
algebraic bound in Lemma~\ref{lem:good23-product} excludes this wall-clock
event and therefore is not a bound on the total deployed Abort probability.
Applying the hiding theorem to that deadline requires either a
witness-independent readiness argument or a padded scheduling mechanism.

The privacy result also does not cover fee-payer linkage, wallet or relayer
identity, transaction origin, the account graph beyond the public statement,
or network-layer metadata.  Those are deployment privacy questions rather
than witness-hiding properties of the proof transcript.

The accepting instruction verifies and mutates state in one transaction only
after a proof account has been created, funded, uploaded, checked, and
finalized.  Those preceding transactions, their latency, fees, rent, and
storage are part of system operation but not part of the one-transaction
claim.  Finalization is an application-level state enforced by the frozen
program: a program upgrade, account close or reallocation path, or governance
change can alter that enforcement.  A deployed program should therefore be
identified together with its ProgramData account, code hash, remaining
capacity, finality status, and upgrade-authority state.

The hiding theorems are conditional on the complete affine-image and
rank-coverage premise in Assumption~\ref{lem:complete-affine-image}.  The
\(\mathsf{Good}_{23}\) checker supplies release-bound evidence by reconstructing
and eliminating the relevant finite-field matrices from the frozen layout.
The fingerprints and regression tests detect drift, but no independent
verifier of the reconstruction and implementation binding is presently
published; the executable evidence is not a formal proof-assistant
development or an external security audit.
Similarly, successful local or devnet execution establishes behavior in the
measured environment; it does not by itself establish mainnet capacity,
operational readiness, historical priority, or resistance to denial-of-
service and transaction-ordering attacks.
